\section{卡特兰数的生成函数：从组合方程到通项}


第 12 讲已经用反射原理得到卡特兰数 $C_n=\dfrac{1}{n+1}\dbinom{2n}{n}$。这一节用生成函数\textbf{重新推出}这个公式：先由组合分解写出 $C(x)$ 满足的代数方程，再解方程、用广义二项式定理展开。这条路线是"生成函数解组合方程"的典型例子，也是下一节 $3\times n$ 多米诺覆盖方法的雏形。

\textbf{本节用到的记号：} $C_n$ 为第 $n$ 个卡特兰数（$C_0=1$：$n$ 对合法括号、Dyck 路、$(n+2)$ 边形三角剖分等）；$C(x)=\sum_{n\ge0}C_nx^n$；$\binom{\alpha}{k}$ 为广义二项式系数（$\alpha$ 可为任意实数）。

\subsection{组合分解给出代数方程}


\textbf{回顾卡特兰数的递推.} $C_0=1$，且对 $n\ge1$，

\[
C_n=\sum_{k=0}^{n-1}C_k\,C_{n-1-k}.
\]

\textbf{用生成函数翻译.} 设 $C(x)=\sum_{n\ge0}C_nx^n$。右边是一个卷积：$C(x)\cdot C(x)$ 的 $x^{n-1}$ 项系数正好是 $\sum_{k=0}^{n-1}C_kC_{n-1-k}=C_n$。因此对 $n\ge1$，$C_n=[x^{n-1}]C(x)^2$。乘以 $x^n$ 并对 $n\ge1$ 求和：

\[
C(x)-C_0=x\,C(x)^2.
\]

代入 $C_0=1$，得到\textbf{卡特兰生成函数方程}：

\[
C(x)=1+xC(x)^2.
\]

\textbf{组合意义.} 一个卡特兰对象（比如合法括号序列）可以唯一分解为

\[
\text{（一对括号包住的内部）}\ \text{后接}\ \text{（剩余部分）},
\]

内部是任意卡特兰对象（权 $C(x)$），外层括号与"后接"各贡献一个 $x$ 的"大小"，所以总生成函数满足 $C=1+xC^2$。这里的"$1$"对应空对象。$3\times n$ 覆盖、二叉树、三角剖分都可以做类似的分解，只是分解方式不同、方程不同。

\subsection{解二次方程}


把 $C=1+xC^2$ 整理成关于 $C$ 的二次方程：

\[
xC^2-C+1=0.
\]

由求根公式，

\[
C(x)=\frac{1\pm\sqrt{1-4x}}{2x}.
\]

\textbf{选哪个分支？} 生成函数在 $x=0$ 附近应等于 $C_0=1$。当 $x\to0$ 时，$\sqrt{1-4x}\to1$，分子 $1-\sqrt{1-4x}\to0$，与分母 $2x\to0$ 的比趋于有限值；而取"$+$"号时分子趋于 $2$，比值发散。因此取\textbf{减号}：

\[
C(x)=\frac{1-\sqrt{1-4x}}{2x}.
\]

这个式子本身已经是一个完整的答案：要算 $C_n$，只需把它展开。

\subsection{用广义二项式定理展开}


\textbf{广义二项式定理.} 对任意实数 $\alpha$ 与 $|x|$ 足够小（形式幂级数意义下则无需收敛），

\[
(1+x)^\alpha=\sum_{k\ge0}\binom{\alpha}{k}x^k,\qquad
\binom{\alpha}{k}=\frac{\alpha(\alpha-1)\cdots(\alpha-k+1)}{k!}.
\]

取 $\alpha=\frac12$、把 $x$ 换成 $-4x$：

\[
\sqrt{1-4x}=\sum_{k\ge0}\binom{1/2}{k}(-4x)^k.
\]

\textbf{求 $C_n$.} 由 $C(x)=\dfrac{1-\sqrt{1-4x}}{2x}$，

\[
C_n=[x^n]C(x)=\frac12[x^{n+1}]\left(1-\sqrt{1-4x}\right)
=-\frac12[x^{n+1}]\sqrt{1-4x}.
\]

（$[x^{n+1}]1=0$ 对 $n\ge0$。）代入展开式：

\[
C_n=-\frac12\binom{1/2}{n+1}(-4)^{n+1}.
\]

下面化简。先算 $\binom{1/2}{n+1}$：

\[
\binom{1/2}{n+1}
=\frac{\frac12\left(-\frac12\right)\left(-\frac32\right)\cdots\left(\frac12-n\right)}{(n+1)!}
=\frac{(-1)^n\,1\cdot3\cdot5\cdots(2n-1)}{2^{n+1}(n+1)!}.
\]

而 $1\cdot3\cdots(2n-1)=\dfrac{(2n)!}{2^nn!}$，所以

\[
\binom{1/2}{n+1}=\frac{(-1)^n(2n)!}{2^{2n+1}\,n!\,(n+1)!}.
\]

代回：

\[
C_n=-\frac12\cdot\frac{(-1)^n(2n)!}{2^{2n+1}n!(n+1)!}\cdot(-1)^{n+1}2^{2n+2}
=\frac{(2n)!}{n!(n+1)!}
=\frac{1}{n+1}\binom{2n}{n}.
\]

这就是卡特兰数的通项公式。\textbf{注意最后一步的因子}：$-\frac12\cdot(-1)^{2n+1}\cdot2^{2n+2}/2^{2n+1}=-\frac12\cdot(-1)\cdot2=1$，所以符号与 $2$ 的幂都恰好抵消——这是推导里最容易算错的地方。

\textbf{验证前几项.} 公式给出

\[
C_0,C_1,C_2,C_3,C_4,C_5=1,1,2,5,14,42,
\]

与递推一致。

\subsection{为什么这仍然是"普通"生成函数}


$C(x)=1+xC^2$ 里的乘法是普通卷积：把两个卡特兰对象按"大小相加"拼接。因此 $C(x)$ 是\textbf{普通生成函数}。若改用指数生成函数 $\hat C(x)=\sum C_nx^n/n!$，递推就不再做普通卷积，方程也不会这么简单（下一份讲解稿会专门对比这两种生成函数）。

\textbf{方法总结（组合方程 → 通项）.}

1. 找到组合对象的递归分解（"第一个块"/"最外层"/"根"等），写出 $C_n$ 与更小 $C_k$ 的关系；
2. 设生成函数 $C(x)$，把递推翻译成含 $C(x)$ 的代数方程（乘积对应卷积、移位对应乘 $x$）；
3. 解方程（二次方程用求根公式，注意用初值选分支）；
4. 用广义二项式定理/部分分式展开，读出系数。

这套流程对 $3\times n$ 覆盖、二叉树、平面树等都适用，只是第 1 步的分解不同。

\subsection{本节的提醒}


\begin{itemize}
\item \textbf{分支选择用初值，不用数值大小。} 选 $\frac{1-\sqrt{1-4x}}{2x}$ 是因为 $C(0)=1$；不要因为"$-\sqrt{}$ 看起来小"就选它。
\item \textbf{广义二项式系数要算对符号。} $\binom{1/2}{k}$ 的符号是 $(-1)^{k-1}$（$k\ge1$），与 $(-4)^k$ 的符号相乘后要仔细化简。
\item \textbf{$C=1+xC^2$ 中的 $x$ 记"一对括号"（或一个内部结点）的大小}；若换一个变量记法，方程形式会变，但方法不变。
\item \textbf{不要把它和"指数生成函数"的方程混淆。} 指数生成函数下，乘法的系数会多出二项式系数；直接套 $C=1+xC^2$ 会得到错误的数列。
\item \textbf{手写稿底部的方程字迹像 $x^2C^2$}，按上面的卷积推导应为 $xC^2$（一个 $x$ 记录"一对括号/一个内部结点"）；整理版统一写 $C=1+xC^2$。
\end{itemize}
